\worksheettitle
  {M05-R. Читаем границу}
  {\modulemark{M05-R} Коррекционная карточка}

\studentnameblock

\begin{warningbox}[Назовите ошибку]
  Вершины --- точки, а звенья и стороны --- отрезки между соседними
  вершинами. Тип фигуры определяют по реально проведённой границе.
\end{warningbox}

\begin{practicebox}[Шаг 1. Посчитайте переходы]
  У ломаной $A-B-C-D$:
  \[
    \text{вершин } \answerblank[12mm],\qquad
    \text{звеньев } \answerblank[12mm].
  \]
  Запишите звенья: \answerblank[75mm].
\end{practicebox}

\begin{practicebox}[Шаг 2. Проверьте замкнутость]
  Ломаная $K-L-M-N-K$ возвращается в начальную точку, поэтому она
  \answerblank[30mm]. Ломаная $P-Q-R-S$ не возвращается, поэтому она
  \answerblank[30mm].
\end{practicebox}

\begin{practicebox}[Шаг 3. Назовите многоугольник]
  Граница $A-B-C-D-A$ состоит из \answerblank[15mm] сторон и образует
  \answerblank[35mm]. Его стороны:
  \answerblank[80mm].
\end{practicebox}

\begin{practicebox}[Шаг 4. Исправьте]
  Ученик сказал: «У незамкнутой ломаной $M-N-P-R$ четыре звена».

  Причина: \answerblank[95mm]

  Исправление: \answerblank[95mm]
\end{practicebox}

\begin{checkpointbox}[Повторная проверка]
  Для границы $E-F-G-H-K-E$ укажите число вершин и сторон, название
  многоугольника и перечислите все стороны. \writinglines{3}
\end{checkpointbox}
